\section{Background}


\subsection{Post-Quantum Cryptography}

NIST standardized three post-quantum algorithms in 2022 \cite{nist2022}:

\begin{table}[h]
\centering
\begin{tabular}{llll}
\toprule
\textbf{Algorithm} & \textbf{Type} & \textbf{Sig Size} & \textbf{Security} \\
\midrule
ML-DSA (FIPS 204, formerly CRYSTALS-Dilithium) & Lattice & 3,293 bytes & NIST Level 3 \\
SLH-DSA (FIPS 205, formerly SPHINCS+) & Hash-based & 17,088 bytes & NIST Level 3 \\
FALCON & Lattice & 1,280 bytes & NIST Level 5 \\
\bottomrule
\end{tabular}
\caption{NIST post-quantum signature schemes}
\end{table}

**Lux choice**: ML-DSA for balance of size and performance, SLH-DSA for stateless long-term security.

\subsection{Lattice-Based Cryptography}

ML-DSA security relies on \textbf{Module Learning With Errors (MLWE)}:

\begin{equation}
\mathbf{b} = \mathbf{A} \cdot \mathbf{s} + \mathbf{e} \pmod{q}
\end{equation}

where:
\begin{itemize}
    \item $\mathbf{A} \in \mathbb{Z}_q^{k \times \ell}$ is public matrix
    \item $\mathbf{s} \in \mathbb{Z}_q^\ell$ is secret key
    \item $\mathbf{e}$ is small error vector
    \item $\mathbf{b}$ is public key
\end{itemize}

**Hardness**: No known quantum algorithm solves MLWE faster than classical lattice reduction ($2^{128}$ security).

\subsection{Lux Multi-Consensus}

Lux employs a family of consensus protocols:

\begin{itemize}
    \item \textbf{Wave}: Linear chain, single-slot finality
    \item \textbf{Lux DAG}: DAG-based, parallel execution
    \item \textbf{Quasar}: Optimistic consensus with fraud proofs
\end{itemize}

**Key property**: Metastability-based consensus (unlike Nakamoto or BFT).

